\section{Heisenberg factors with fully deleted binary labels}
\label{sec:heisenberg-deleted-labels}

\paragraph{Author freeze and origin.}
D-TD22--23, reserved question Q2el, is a derivation by
\texttt{/root}. Its elementary binary-rank estimate supplies a
separate actual-label case of the residual Heisenberg problem.
The complete author proof, dependency-interface, checker and first-output
readings and the complete independent whole-proof review pass.
No classification is used to assert that an arbitrary label module
has this form. The input trace formula is the accepted D-TD05,
and the actual-source obstruction is HD-FT27.

Let $n\ge5$, let $M=\mathbf F_2^n$ be the natural permutation
module, and put
\[
 A=\{(z_1,\ldots,z_n):\textstyle\sum_i z_i=0\},\qquad
 T=\mathbf F_2(1,\ldots,1),\qquad
 Q_n=A/(A\cap T).
 \tag{TDB.1}
\]
The dot product on $M$ induces a nondegenerate alternating form
on $Q_n$. Indeed $A^\perp=T$, and for $v\in A$ its self-pairing
is $\sum_i v_i^2=\sum_i v_i=0$. Thus
\[
 d_n:=\dim_{\mathbf F_2}Q_n=
 \begin{cases}n-1&n\text{ odd},\\n-2&n\text{ even},\end{cases}
 \qquad d_n\text{ is even}.
 \tag{TDB.2}
\]
Write $f(y)$ for the number of fixed natural letters of $y\in S_n$,
and $\operatorname{ind}(y)$ for $n$ minus its number of cycles,
including fixed cycles.

\subsection*{D-TD22. Binary deleted ranks and cubic Heisenberg traces}

\begin{lemma}[D-TD22]
For $y\in S_n$, the rank on $Q_n$ satisfies
\[
 \operatorname{rank}_{Q_n}(1-y)\ge\operatorname{ind}(y)-2,
 \qquad y\ne1\Longrightarrow
                \operatorname{rank}_{Q_n}(1-y)\ge1.
 \tag{TDB.3}
\]
Suppose $Q_n$ is an actual elementary abelian projective subgroup
on an irreducible complex Heisenberg space $W$, with nondegenerate
scalar commutator form and $a=\dim W=2^{d_n/2}$.
For any finite-order normalizing lift $B$ inducing the natural
action of $y$ on these labels, one has
\[
 \frac{|\operatorname{tr}B|}{a}
 \le 2^{-\operatorname{rank}_{Q_n}(1-y)/2}.
 \tag{TDB.4}
\]
If $n\ge4096$, then $a>n^3$ and, for every such lift,
\[
 \frac{|\operatorname{tr}B|}{a}
 \le\left(\frac{\max\{f(y),1\}}n\right)^3,
 \qquad y\ne1\Longrightarrow
 \frac{|\operatorname{tr}B|}{a}\le\frac1{\sqrt2}<\frac78.
 \tag{TDB.5}
\]
These bounds include every power of every finite-order lift;
no equality of its order with the natural permutation order is assumed.
\end{lemma}

\begin{proof}
On $M$, the kernel of $1-y$ consists of vectors constant on each
cycle, in every characteristic. Its rank is therefore exactly
$\operatorname{ind}(y)$. Restricting a linear map to an invariant
subspace of codimension one loses at most one rank. Quotienting
by an invariant subspace of dimension at most one loses at most
one more. This proves the first inequality of (TDB.3).

We verify faithfulness on $Q_n$ to cover small permutation indices
where that first estimate alone could be zero. If $y$ acts trivially
on $Q_n$, then for every pair of distinct letters $i,j$,
\[
 e_{y(i)}+e_{y(j)}+e_i+e_j\in A\cap T.
 \tag{TDB.6}
\]
Its support has at most four letters. Since $n\ge5$, it cannot
be the all-one vector, so it is zero. Thus $y$ fixes every
two-element subset. The intersection of the two subsets
$\{i,j\}$ and $\{i,k\}$, with $i,j,k$ distinct, shows that it
fixes each letter. Hence $y=1$, proving the second inequality.

The complete intrinsic normalizer trace formula of D-TD05 says
that $|\operatorname{tr}B|^2$ is either zero or
$2^{\dim\ker(1-y)}$ for a unitary normalizing lift.
The same accepted theorem identifies the projective centralizer
of the label group with that finite group itself. Its full projective
normalizer is therefore finite, since its quotient embeds in the
finite symplectic group. Any finite-order lift under consideration
can thus be represented by a unitary lift with the same trace modulus.
More explicitly, the
full inverse image of that finite projective image in
$\operatorname{SL}(W)$ is finite; averaging gives a common invariant
Hermitian form. A finite-order scalar multiplying such a lift has
absolute value one. Dividing the trace formula by $a^2=2^{d_n}$
gives (TDB.4). Equivalently, the formula applies directly to each
finite-order normalizer after choosing its invariant Hermitian form.
No quadratic refinement or split normalizer extension is used.

Here is a universal elementary exponential margin. For $n\ge4096$
choose the integer $q\ge12$ with $2^q\le n<2^{q+1}$.
At $q=12$ one has $2^{q-6}-1>3(q+1)$, and increasing $q$ by one
increases the left side by $2^{q-6}\ge64$ and the right side by
three. Therefore
\[
 2^{n/64-1}>2^{3(q+1)}>n^3,
 \qquad
 a\ge2^{(n-2)/2}\ge2^{n/64-1}>n^3.
 \tag{TDB.7}
\]
This is an induction for all $q\ge12$, not an extrapolation
from a finite grid.

Put $k=n-f(y)$. For $y\ne1$, (TDB.3)--(TDB.4) always give
the uniform ratio at most $1/\sqrt2$. If $k\le n/16$, then
\[
 \frac1{\sqrt2}\le\left(\frac{15}{16}\right)^3
                     \le\left(\frac{f(y)}n\right)^3,
 \tag{TDB.8}
\]
where the first inequality is the exact integer comparison
$2\cdot15^6\ge16^6$.
If $k>n/16$, every nontrivial cycle of length $l$ contributes
$l-1\ge l/2$ to the natural index. Thus
\[
 \operatorname{rank}_{Q_n}(1-y)\ge\operatorname{ind}(y)-2
          \ge k/2-2>n/32-2,
 \qquad
 \frac{|\operatorname{tr}B|}{a}
       <2^{1-n/64}<\frac1{n^3}.
 \tag{TDB.9}
\]
This is at most the desired cubic expression, also when $f(y)=0$.
When $y=1$, finite order gives the needed ratio at most one.
Finally $1/\sqrt2<7/8$ follows by squaring, since $1/2<49/64$.
Every power of a finite-order normalizer is another such lift,
with natural action $y^j$; the argument therefore proves precisely
the full-power assertion as well.
\end{proof}

\subsection*{D-TD23. Exclusion with an arbitrary whole partner}

\begin{theorem}[D-TD23]
Let $J$ be an actual faithful finite complex coefficient image
on $V$, and suppose it preserves the whole matrix algebra of
a Heisenberg factor $W$ with an actual nondegenerate label group
isomorphic, as a symplectic module with its entire conjugation image,
to $Q_n$ with $P=A_n$ or $S_n$ in the action (TDB.1).
Retain the entire external partner:
\[
 V=W\otimes R,\qquad b=\dim R\ge1,\qquad
 e=ab,\qquad x:J\longrightarrow P\text{ onto}.
 \tag{TDB.10}
\]
For $n\ge10^6$, such a coefficient image cannot be the coefficient
representation of an actual connected sphere cover with moving-rank
sum $2e+6$. It therefore cannot be an actual rank-nine sixfold
inducing source in the established high-degree source interface.
No maximality, primitivity, coefficient field, or support floor
on the entire partner $R$ is required.
\end{theorem}

\begin{proof}
Stability of the full factor algebra gives a genuine finite
projective image $P_W$ of $J$ on $W$. Its conjugation action on
the actual elementary abelian label group gives the displayed
global surjection $x$. This is the stated entire action, not an
alternating quotient of an unspecified different subgroup.

For $g\in J$, choose a lift $B_g$ in the finite inverse image
of $P_W$ in $\operatorname{SL}(W)$ and absorb the reciprocal scalar
in the other factor, giving $V(g)=B_g\otimes C_g$ exactly.
Let $L$ be the order of $V(g)$. Since this tensor has $L$th power
one, each of its factor $L$th powers is scalar. For example,
conjugating the full first factor algebra shows $B_g^L=cI$,
and the tensor identity gives $C_g^L=c^{-1}I$.
The scalar $c$ is a root of unity because $B_g$ is finite order.
Thus $C_g$ is also finite order. For every integer $j$ one obtains
\[
 \frac{|\operatorname{tr}_V(g^j)|}{e}
 =\frac{|\operatorname{tr}(B_g^j)|}{a}
     \frac{|\operatorname{tr}(C_g^j)|}{b}
 \le\frac{|\operatorname{tr}(B_g^j)|}{a}.
 \tag{TDB.11}
\]
No coherent ordinary linear action of $J$ on either chosen tensor
lift has been asserted, and no kernel power has been dropped.
The natural action of $B_g^j$ is $x(g)^j$.

D-TD22 gives both whole-coefficient trace hypotheses of HD-FT27,
and gives $e\ge a>n^3$. The actual quotient $x$ and connected
sphere cover supply the remaining hypotheses. HD-FT27 contradicts
the moving-rank sum $2e+6$ using its full original-order Fourier
and actual cube-quotient genus argument. The established inducing
source Euler/Hurwitz interface, D-BI-18 at (BI.65)--(BI.66),
forces that same sphere setting for a rank-nine source at these
degrees, so the stated application follows.
\end{proof}

\paragraph{Boundary.}
This is an explicit represented-label exclusion. It does not replace
an arbitrary irreducible subgroup of $\operatorname{Sp}(2m,p)$
by the full group, lift an alternating quotient to a subgroup,
classify arbitrary homogeneous modular kernels, or assert that
the label module is fully deleted from character values alone.
Any classification application must separately prove the exact
module and actual-image hypotheses of D-TD23. No new source,
polarization comparison or effective detector is constructed.
